\documentclass[11pt]{article}
\usepackage{fontspec,amsmath,amssymb,geometry}
\setmainfont{Times New Roman}\geometry{margin=0.85in}
\newcommand{\solutionquestion}[1]{\par\leftskip=0pt\section*{Q#1}\leftskip=1em}
\newcommand{\solutionpart}[1]{\par\smallskip\leftskip=1em\noindent\hangindent=2em\hangafter=1\makebox[1.5em][l]{\textbf{(#1)}}\hspace{0.5em}\ignorespaces}
\newcommand{\solutioncontinuation}{\par\leftskip=3em\noindent\ignorespaces}
\title{Determinants}\author{T.T.}\date{}
\begin{document}\maketitle

\solutionquestion{1}
Apply
\[
R_2\leftarrow R_2-2R_1,
\qquad
R_3\leftarrow R_3-R_1.
\]
These row additions do not change the determinant, and they give
\[
\begin{bmatrix}
1&2&3\\
0&1&1\\
0&-1&-2
\end{bmatrix}.
\]
Now apply $R_3\leftarrow R_3+R_2$, again without changing the determinant:
\[
U=\begin{bmatrix}
1&2&3\\
0&1&1\\
0&0&-1
\end{bmatrix}.
\]
There were no row exchanges or row scalings. The pivots are $1$, $1$, and $-1$, so
\[
\det A=\det U=1\cdot1\cdot(-1)=-1.
\]

\solutionquestion{2}
\solutionpart{i}The product rule and transpose rule give
\[
\det(AB)=(\det A)(\det B)=(-2)(3)=-6,
\qquad
\det({A}^{T})=\det A=-2.
\]

\solutionpart{ii}Since $AA^{-1}=I$,
\[
\det A\,\det(A^{-1})=\det I=1,
\]
so
\[
\det(A^{-1})=-\frac12.
\]
Also,
\[
\det({A}^{2})={\det A}^{2}=4,
\qquad
\det(2A)=2^{4}\det A=-32.
\]

\solutionquestion{3}
\solutionpart{i}The cofactors are
\[
C=\begin{bmatrix}
-5&-1&4\\
-1&2&-8\\
2&-4&5
\end{bmatrix}.
\]
For example,
\[
C_{11}=\det\begin{bmatrix}3&2\\4&1\end{bmatrix}=-5,
\qquad
C_{12}=-\det\begin{bmatrix}1&2\\0&1\end{bmatrix}=-1.
\]
The other entries follow in the same way.

\solutionpart{ii}Expanding along the first row gives
\[
\begin{aligned}
\det A
&=2C_{11}+1C_{12}+0C_{13}\\
&=2(-5)+(-1)=-11.
\end{aligned}
\]

\solutionpart{iii}Using the cofactor matrix,
\[
A{C}^{T}
=\begin{bmatrix}2&1&0\\1&3&2\\0&4&1\end{bmatrix}
\begin{bmatrix}-5&-1&2\\-1&2&-4\\4&-8&5\end{bmatrix}
=\begin{bmatrix}-11&0&0\\0&-11&0\\0&0&-11\end{bmatrix}
=-11I=(\det A)I.
\]

\solutionquestion{4}
\solutionpart{i}Subtract the first row from the second and third rows:
\[
\begin{bmatrix}1&1&1\\0&t-1&1\\0&1&t-1\end{bmatrix}.
\]
Expanding along the first column gives
\[
\det A(t)={(t-1)}^{2}-1=t(t-2).
\]

\solutionpart{ii}The matrix is invertible exactly when its determinant is nonzero. Therefore
\[
A(t)\text{ is invertible if and only if }t\ne0\text{ and }t\ne2.
\]

\solutionpart{iii}For $t=0$, the equations $A(0)x=0$ reduce to
\[
x+y+z=0,\qquad x+2z=0,\qquad x+2y=0.
\]
Thus $(-2,1,1)$ is a nonzero null vector. For $t=2$, the last two rows agree, and
\[
(0,1,-1)
\]
is a nonzero null vector because each row has zero dot product with it.

\solutionquestion{5}
Write the system as $Ax=b$, where
\[
A=\begin{bmatrix}2&1&0\\1&2&1\\0&1&2\end{bmatrix},
\qquad
b=\begin{bmatrix}1\\2\\3\end{bmatrix}.
\]
The determinant is
\[
D=\det A
=2\det\begin{bmatrix}2&1\\1&2\end{bmatrix}
-\det\begin{bmatrix}1&1\\0&2\end{bmatrix}
=2(3)-2=4.
\]
Replacing the first column by $b$ gives
\[
D_1=\det\begin{bmatrix}1&1&0\\2&2&1\\3&1&2\end{bmatrix}=2.
\]
Replacing the second column gives
\[
D_2=\det\begin{bmatrix}2&1&0\\1&2&1\\0&3&2\end{bmatrix}=0.
\]
Replacing the third column gives
\[
D_3=\det\begin{bmatrix}2&1&1\\1&2&2\\0&1&3\end{bmatrix}=6.
\]
By Cramer's rule,
\[
x=\frac{D_1}{D}=\frac12,
\qquad
y=\frac{D_2}{D}=0,
\qquad
z=\frac{D_3}{D}=\frac32.
\]

\solutionquestion{6}
The cofactor matrix of $A$ is
\[
C=\begin{bmatrix}
0&-2&2\\
-2&1&-1\\
2&-1&-3
\end{bmatrix}.
\]
Expanding along the first row gives
\[
\det A=1(0)+2(-2)+0=-4.
\]
The cofactor formula therefore gives
\[
A^{-1}=\frac{1}{\det A}{C}^{T}
=-\frac14\begin{bmatrix}
0&-2&2\\
-2&1&-1\\
2&-1&-3
\end{bmatrix}.
\]
The cofactor identity yields
\[
A{C}^{T}=-4I,
\]
and hence
\[
AA^{-1}=A\left(-\frac14{C}^{T}\right)=I.
\]

\solutionquestion{7}
\solutionpart{i}Put the vectors into the columns of a matrix. The parallelogram area is
\[
\left|\det\begin{bmatrix}3&1\\2&4\end{bmatrix}\right|
=|12-2|=10.
\]

\solutionpart{ii}The volume is the absolute value of the determinant whose columns are $u$, $v$, and $w$:
\[
\begin{aligned}
\operatorname{Vol}
&=\left|\det\begin{bmatrix}1&1&0\\1&0&1\\0&1&1\end{bmatrix}\right|\\
&=|-2|=2.
\end{aligned}
\]

\solutionquestion{8}
\solutionpart{i}The cofactor formula is
\[
A^{-1}=\frac{1}{\det A}{C}^{T}.
\]
Because $A$ has integer entries, every minor and every cofactor is an integer. If $\det A=1$ or $\det A=-1$, division by $\det A$ changes only the signs of the cofactors. Therefore every entry of $A^{-1}$ is an integer.

\solutionpart{ii}If $A$ and $A^{-1}$ both have integer entries, then their determinants are integers. Since
\[
\det A\,\det(A^{-1})=\det I=1,
\]
two integers multiply to $1$. The only possibilities are
\[
\det A=1,\quad\det(A^{-1})=1,
\]
or
\[
\det A=-1,\quad\det(A^{-1})=-1.
\]
Thus $\det A=1$ or $\det A=-1$.

\end{document}
